% ID: SEL-D0004-OA-E04
% ESTADO: validada. S005, 2026-11-07; tres pasadas del mismo autor y revisión matemática, didáctica y visual.
\textbf{Modelo.}\par\nopagebreak[4]
Se estima la proporción de quienes \emph{no} probaron el sabor: $\widehat p=20/125=0.16$, no 0.84. Bajo muestreo aleatorio representativo e independencia, el intervalo normal aproximado es
\[
\widehat p\pm z_{1-\alpha/2}\sqrt{\frac{\widehat p(1-\widehat p)}n}.
\]
Los recuentos 20 y 105 permiten la aproximación usual. La representatividad de la encuesta es una condición del procedimiento, no algo que demuestre el tamaño muestral.

\textbf{a) Confianza del 97\,\%.}\par\nopagebreak[4]
$\alpha=0.03$ y $1-\alpha/2=0.985$, de donde $z\approx2.170090378$ (tabla: 2.17). El error es
\[
E=z\sqrt{\frac{0.16(0.84)}{125}}\approx0.071157790.
\]
\[
\boxed{IC_{\text{97\%}}\approx(0.088842210,\ 0.231157790).}
\]
Equivale aproximadamente a entre 8.88\,\% y 23.12\,\%. Es un intervalo para la proporción poblacional; la confianza describe la cobertura del método en encuestas repetidas bajo las hipótesis indicadas.

\textbf{b) Nueva confianza del 94\,\% y error estricto.}\par\nopagebreak[4]
Ahora $1-\alpha/2=0.97$ y $z\approx1.880793608$, manteniendo la proporción piloto 0.16. El error inferior al 5\,\% significa $E<0.05$:
\[
z\sqrt{\frac{0.16(0.84)}n}<0.05
\iff n>\frac{z^2(0.16)(0.84)}{0.05^2}\approx190.169796.
\]
Así, \boxed{n=191\text{ personas}} es el tamaño total mínimo. No es el número de encuestados adicionales. Se verifica el límite calculando el error para 190 y 191; el primero supera 0.05 y el segundo es menor. Usar el cuantil de tabla 1.88 también da 191. No se sustituye 0.16 por 0.5 porque el problema exige mantener la proporción anterior.
